\begin{abstract}
Pass-end prediction of grinding surface roughness (\(R_a\)) from signals
acquired in process may support
earlier detection of surface-quality deviations, but small condition-structured
datasets make model selection, validation, and interpretation difficult. This
study benchmarks 27 model variants across 38 input configurations using
strict leave-one-condition-out cross-validation over 16 grinding conditions.
Inputs include process parameters, scalar signal features, full-resolution dB
spectrograms, per-sample-standardised dB (dB-z) spectrograms, log-mel
spectrograms, and two-dimensional wavelet-scattering descriptors.

The leading observed point-estimate group comprises random forests on fused
acoustic-emission and vibration dB-z and log-mel descriptors. The
reproducible current fixed dB-z pipeline gives \(0.0210~\upmu\)m, while
cache-matched inner-grouped reruns give \(0.0207~\upmu\)m for dB-z and
\(0.0206~\upmu\)m for log-mel. Their mean errors are nearly tied, but the
nested maximum fold errors are 0.1402 and 0.2061~\(\upmu\)m, respectively,
so dB-z has better observed worst-condition robustness. Historical canonical ranking artifacts
(0.0200 and 0.0201~\(\upmu\)m) are retained for provenance and the planned
post-selection comparison, not as evidence of a unique winner. For the
current fixed dB-z model, the median fold MAE is \(0.0093~\upmu\)m but the
maximum is \(0.1478~\upmu\)m on Condition~7, exposing substantial unseen-condition
tail risk. These differences are also close to the available measurement
uncertainty scale. In the planned post-selection family of 14
Holm--Bonferroni-corrected Wilcoxon tests, only the comparison with the
implementation-specific historical shallow MLP remains significant after
correction; a corrected whole-condition-validation MLP sensitivity remains
worse but is not part of that historical family. Vibration alone nearly
matches the fused RF, so the observed mean-level contribution of AE is small.

TreeSHAP, ridge coefficients, and Grad-CAM identify model-associated
frequency regions, but sensor-response limitations preclude causal
interpretation of high-frequency AE attributions. For a separately regenerated
ResNetVibCNN checkpoint set with deterministic MAE of \(0.0443~\upmu\)m,
nominal MC-dropout epistemic-band coverage is 50.2\%; out-of-bag-scaled intervals for a separate 15-condition RF uncertainty variant also
under-cover difficult conditions, including
Condition~7. The results support auditable random forests with dB-z or
log-mel pass-level spectral descriptors as strong observed baselines in this
dataset, rather than general superiority over deep learning. The comparison
is post-selection and non-nested, and full-pass feature aggregation does not
establish within-pass or closed-loop readiness.
\end{abstract}
